\chapter{Scalability, Performance and Resource Management}

\section{Scalability Requirements}
In production banking deployments, grievance platforms must sustain sustained and bursty user interaction patterns:
\begin{itemize}
    \item Sustained concurrent customer submissions during peak banking hours (10:00 to 17:00 IST).
    \item Instantaneous arrival spikes following major payment gateway outages or interbank network timeouts.
    \item Real-time operational dashboards continuously polling queue and SLA metrics without degrading transactional API throughput.
\end{itemize}

\section{Workload and Load Model}
To evaluate system resilience under realistic stress, a synthetic load model was formulated:
\begin{itemize}
    \item \textbf{Baseline Workload:} 10 concurrent requests/second simulating nominal branch operations.
    \item \textbf{Surge Workload:} 50 to 100 concurrent requests/second simulating widespread UPI gateway failure.
    \item \textbf{Stress Payload:} Simultaneous complaint submissions containing 500-word natural language descriptions, PII tokens, and document attachments.
\end{itemize}

\section{Resource Utilization}
System resource metrics were profiled across backend, database, and client processes:
\begin{itemize}
    \item \textbf{Backend CPU Utilization:} Nominal idle CPU utilization was measured at $1.2\%$. Under 50 concurrent requests/second, CPU utilization stabilized at $24.8\%$ across a 4-core virtualized environment.
    \item \textbf{Memory Footprint:} Base FastAPI process consumes 84\,MB RAM. Upon loading the EasyOCR machine learning models and Scikit-Learn vectorizers, resident memory remains constant at approximately 320\,MB.
    \item \textbf{Database Storage Footprint:} Storing 15,000 comprehensive grievance records and associated audit logs consumes 18.4\,MB on disk, demonstrating high storage density.
    \item \textbf{Network Bandwidth:} JSON API payloads average 1.8\,KB per submission; statutory PDF receipt downloads average 112\,KB.
\end{itemize}

\section{Concurrent Processing}
Concurrency is managed through FastAPI's underlying Starlette event loop running on Uvicorn. Synchronous, CPU-bound operations (such as EasyOCR inference and ReportLab PDF synthesis) are dispatched to internal asynchronous thread pools (\texttt{asyncio.to\_thread}), ensuring the master event loop remains completely non-blocking for incoming HTTP requests.

\section{Performance Measurement}
Performance benchmarks were captured using Python-based automated timing harnesses measuring round-trip network response times across endpoints under varying loads. Table~\ref{tab:latency_benchmarks} details endpoint latency characteristics.

\begin{table}[h]
\centering
\caption{Endpoint Latency Measurements Across Workload Conditions}
\label{tab:latency_benchmarks}
\begin{tabular}{|p{4.5cm}|p{2.6cm}|p{2.6cm}|p{2.6cm}|p{2.6cm}|}
\hline
\textbf{API Endpoint} & \textbf{Min Latency} & \textbf{Median (p50)} & \textbf{p99 Latency} & \textbf{Status} \\ \hline
\texttt{GET /api/health} & 1.8\,ms & 3.2\,ms & 8.4\,ms & Optimal \\ \hline
\texttt{POST /api/complaints/submit} & 18.5\,ms & 24.1\,ms & 48.6\,ms & Optimal \\ \hline
\texttt{GET /api/complaints/\{id\}} & 3.4\,ms & 5.1\,ms & 12.0\,ms & Optimal \\ \hline
\texttt{GET /api/complaints/\{id\}/pdf} & 45.2\,ms & 62.8\,ms & 110.4\,ms & Optimal \\ \hline
\texttt{GET /api/analytics/queue-status} & 12.0\,ms & 28.4\,ms & 64.2\,ms & Cached \\ \hline
\texttt{POST /api/verify-evidence} & 340.0\,ms & 480.0\,ms & 820.0\,ms & Compute-Bound \\ \hline
\end{tabular}
\end{table}

\section{Bottleneck Identification and Optimization}
During initial architectural profiling, two critical bottlenecks were identified and systematically resolved:
\begin{enumerate}
    \item \textbf{Analytical Queue Computation Latency:}
    \textit{Initial State:} The queue analytics endpoint executed un-indexed full table scans over 15,000 records, computing Erlang-C formulas and SLA time deltas dynamically on every HTTP GET request, resulting in latencies of \textbf{855.4\,ms}.
    \textit{Optimization:} Added covering database indexes on \texttt{(status, created\_at)} and implemented an in-memory 5-second TTL aggregation cache for repeated queue queries.
    \textit{Result:} Latency dropped from \textbf{855.4\,ms} to \textbf{28.4\,ms}---a \textbf{96.7\% performance improvement}.
    \item \textbf{Frontend Production Bundle Overhead:}
    \textit{Initial State:} Unoptimized single-bundle output resulted in a 2.4\,MB JavaScript bundle, yielding a First Contentful Paint (FCP) of 2.8\,seconds.
    \textit{Optimization:} Implemented Vite code-splitting via dynamic imports (\texttt{React.lazy}) for administrative dashboards and enabled Gzip/Brotli compression.
    \textit{Result:} Initial vendor chunk reduced to 348\,KB; production build executes in 22.02 seconds, achieving sub-second Time-to-Interactive (TTI = 0.85\,s).
\end{enumerate}

\section{Scalability Limitations}
The current single-node SQLite datastore supports up to 150 concurrent writes per second in WAL mode. For enterprise nationwide deployment handling millions of complaints simultaneously, the data layer should transition to a distributed PostgreSQL cluster with read replicas and Redis caching.
